% from ega1-10-fr.tex
\oldpage[I]{200}

et un seul, tel que $v^\sharp=(u^\sharp)_{/X'}$. Si on considère les
triplets $(\mathscr{F},X,X')$ ($\mathscr{F}$ étant un
$\mathscr{O}_X$-Module cohérent et $X'$ une partie fermée de $X$) comme
une \emph{catégorie}, les morphismes
$(\mathscr{F},X,X')\to(\mathscr{G},Y,Y')$ consistant en un morphisme de
préschémas $f:X\to Y$ tel que $f(X')\subset Y'$ et en un
$f$-morphisme $u:\mathscr{G}\to\mathscr{F}$, on peut donc dire que
$(X_{/X'},\mathscr{F}_{/X'})$ est un \emph{foncteur} en
$(\mathscr{F},X,X')$, prenant ses valeurs dans la catégorie des couples
$(\mathfrak{Z},\mathscr{H})$ formés d'un préschéma formel localement
noethérien $\mathfrak{Z}$ et d'un
$\mathscr{O}_{\mathfrak{Z}}$-Module $\mathscr{H}$, les morphismes de
cette dernière catégorie consistant en les couples formés d'un morphisme
$g$ de préschémas formels et d'un $g$-morphisme.
\end{env}

\begin{proposition}[10.9.7]
\label{I.10.9.7-fr}
Soient $S$, $X$, $Y$ trois préschémas localement noethériens,
$g:X\to S$, $h:Y\to S$ deux morphismes, $S'$ une partie fermée de $S$,
$X'$ (resp. $Y'$) une partie fermée de $X$ (resp. $Y$) telle que
$g(X')\subset S'$ (resp. $h(Y')\subset S'$)~; soit
$Z=X\times_S Y$~; supposons $Z$ localement noethérien, et soit
$Z'=p^{-1}(X')\cap q^{-1}(Y')$, où $p$ et $q$ sont les projections de
$X\times_S Y$. Dans ces conditions, le complété $Z_{/Z'}$ s'identifie au
produit des $S_{/S'}$-préschémas formels
$(X_{/X'})\times_{S_{/S'}}(Y_{/Y'})$, les morphismes structuraux
s'identifiant à $\widehat g$ et $\widehat h$, et les projections à
$\widehat p$ et $\widehat q$.
\end{proposition}

Il est immédiat que la question est locale pour $S$, $X$ et $Y$, et on
est donc ramené au cas où $S=\operatorname{Spec}(A)$,
$X=\operatorname{Spec}(B)$, $Y=\operatorname{Spec}(C)$,
$S'=V(\mathfrak{J})$, $X'=V(\mathfrak{K})$, $Y'=V(\mathfrak{L})$, où
$\mathfrak{J}$, $\mathfrak{K}$, $\mathfrak{L}$ sont trois idéaux tels
que $\varphi(\mathfrak{J})\subset\mathfrak{K}$ et
$\psi(\mathfrak{J})\subset\mathfrak{L}$, en désignant par $\varphi$ et
$\psi$ les homomorphismes $A\to B$ et $A\to C$ qui correspondent à $g$
et $h$. Alors on sait que $Z=\operatorname{Spec}(B\otimes_A C)$ et que
$Z'=V(\mathfrak{M})$, où $\mathfrak{M}$ est l'idéal
$\operatorname{Im}(\mathfrak{K}\otimes_A C)
+\operatorname{Im}(B\otimes_A\mathfrak{L})$. La conclusion résulte
(\hyperref[I.10.7.2-fr]{10.7.2}) de ce que le produit tensoriel complété
$(\widehat B\otimes_{\widehat A}\widehat C)^\wedge$ (où $\widehat A$,
$\widehat B$, $\widehat C$ sont respectivement les séparés complétés de
$A$, $B$, $C$ pour les topologies $\mathfrak{J}$-, $\mathfrak{K}$- et
$\mathfrak{L}$-préadiques) est le séparé complété du produit tensoriel
$B\otimes_A C$ pour la topologie $\mathfrak{M}$-préadique
(\hyperref[0.7.7.2-fr]{0, 7.7.2}).

On notera en outre que si $T$ est un $S$-préschéma localement
noethérien, $u:T\to X$, $v:T\to Y$ deux $S$-morphismes, $T'$ une partie
fermée de $T$ telle que $u(T')\subset X'$, $v(T')\subset Y'$, alors le
prolongement aux complétés $((u,v)_S)^\wedge$ s'identifie à
$(\widehat u,\widehat v)_{S_{/S'}}$.

\begin{corollary}[10.9.8]
\label{I.10.9.8-fr}
Soient $X$, $Y$ deux $S$-préschémas localement noethériens tels que
$X\times_S Y$ soit localement noethérien~; soient $S'$ une partie fermée
de $S$, $X'$ (resp. $Y'$) une partie fermée de $X$ (resp. $Y$) dont
l'image dans $S$ est contenue dans $S'$. Pour tout $S$-morphisme
$f:X\to Y$ tel que $f(X')\subset Y'$, le morphisme graphe
$\Gamma_{\widehat f}$ s'identifie au prolongement $(\Gamma_f)^\wedge$ du
morphisme graphe de $f$.
\end{corollary}

\begin{corollary}[10.9.9]
\label{I.10.9.9-fr}
Soient $X$, $Y$ deux préschémas localement noethériens, $f:X\to Y$ un
morphisme, $Y'$ une partie fermée de $Y$, $X'=f^{-1}(Y')$. Alors le
préschéma $X_{/X'}$ s'identifie par le diagramme commutatif
\[
  \label{I.10.9.9.1-fr}
  \xymatrix{
    X\ar[d]_f &
    X_{/X'}\ar[l]\ar[d]^{\widehat f}\\
    Y &
    Y_{/Y'}\ar[l]
  }
\]
au produit $X\times_Y(Y_{/Y'})$ de préschémas formels.
\end{corollary}

Il suffit d'appliquer (\hyperref[I.10.9.7-fr]{10.9.7}) en remplaçant
$S$ et $S'$ par $Y$, $X$ et $X'$ par $X$.
